\chapter{Results and Discussion}
\label{chap:results}
% Target ~12 pages -- the 40-point rubric criterion. All tables and
% figures generated by scripts for reproducibility. Include
% uncertainty/error analysis explicitly.

\section{Experiment 1: Evaluative Bias at Scale (H1)}
% Prose written after E1 completes; table auto-generated by
% Code/src/make_tables.py from the experiment logs.
Table~\ref{tab:e1-qwk} reports agreement with ground truth per model,
condition, and cohort.
%% TODO: results prose after E1.
Lorem ipsum dolor sit amet, consectetur adipiscing elit.

% AUTO-GENERATED by make_tables.py -- do not edit by hand.
\begin{table}[htbp]
\centering
\caption[E1: cohort score gap by model]{Experiment E1: mean consensus
score by cohort and the native$-$SEA score gap on proficiency-matched
cohorts, with 95\% bootstrap confidence intervals (10{,}000
resamples).}
\label{tab:e1-qwk}
\small
\begin{tabular}{llccc}
\hline
Model & Condition & Mean score\textsubscript{SEA} &
Mean score\textsubscript{native} & Gap [95\% CI] \\
\hline
Claude Sonnet & Baseline & 4.44 ($n$=100) & 5.62 ($n$=100) & 1.18 [0.86, 1.50] \\
Gemini Pro & Baseline & 4.83 ($n$=100) & 5.00 ($n$=17) & 0.17 [-0.14, 0.49] \\
GPT-4o & Baseline & 5.03 ($n$=100) & 5.83 ($n$=100) & 0.80 [0.60, 1.00] \\
GPT-4o & CAMPS $k{=}3$ & 4.70 ($n$=100) & 5.27 ($n$=100) & 0.57 [0.37, 0.78] \\
\hline
\end{tabular}
\end{table}


\section{Experiment 2: Extended CAMPS Panel (H2, H4)}
Table~\ref{tab:e2-ablation} reports the ablation over panel size and
the meta-adjudicator; the weight-sweep frontier is reported as a
figure generated from the same logs.
%% TODO: results prose after E2.
Lorem ipsum dolor sit amet, consectetur adipiscing elit.

% PLACEHOLDER -- overwritten by make_tables.py after E2 runs.
\begin{table}[htbp]
\centering
\caption[E2: ablation over panel size and adjudicator]{Experiment E2:
ablation over panel size $k$ and the meta-adjudicator. Values pending
completion of E2.}
\label{tab:e2-ablation}
\small
\begin{tabular}{lccc}
\hline
Configuration & QWK\textsubscript{SEA} & QWK\textsubscript{native} & Gap \\
\hline
Baseline ($k{=}1$) & -- & -- & -- \\
CAMPS $k{=}3$ & -- & -- & -- \\
CAMPS $k{=}5$ & -- & -- & -- \\
CAMPS $k{=}5$ + adjudicator & -- & -- & -- \\
\hline
\end{tabular}
\end{table}


\section{Experiment 3: BDS Validation (H3)}
Table~\ref{tab:e3-bds} reports detection performance on the held-out
GRA split.
%% TODO: results prose after E3.
Lorem ipsum dolor sit amet, consectetur adipiscing elit.

% AUTO-GENERATED by make_tables.py -- do not edit by hand.
\begin{table}[htbp]
\centering
\caption[E3: BDS detection performance]{Experiment E3: BDS detection
performance on the held-out GRA split (model: gpt-4o).}
\label{tab:e3-bds}
\small
\begin{tabular}{lc}
\hline
Quantity & Value \\
\hline
Calibrated threshold $\tau$ & 0.400 \\
Sensitivity (held-out, $n$=68) & 0.875 \\
False-positive rate (held-out) & 0.712 \\
Area under ROC curve (held-out) & 0.603 \\
Correlation of BDS with human dispersion & 0.158 \\
\hline
\end{tabular}
\end{table}


\section{Uncertainty and Error Analysis}
% Bootstrap distributions; qualitative analysis of mis-scored essays;
% cases where CAMPS fails; comparison of pilot artifacts (e.g., small-n
% ceiling effects) with at-scale estimates.
%% TODO: Week 8.
Lorem ipsum dolor sit amet, consectetur adipiscing elit.

\section{Threats to Validity}
\label{sec:threats}

\textbf{Construct validity.} The central construct risk is the use of
first-language and country cohorts as proxies for rhetorical
tradition. Rhetorical conventions are distributed across writing
communities, not fixed properties of individuals
\citep{connor2011intercultural}, so cohort-level results must not be
read as claims about any individual writer; this thesis draws
cohort-level conclusions only. A second construct risk concerns the
agents themselves: a prompted ``cultural lens'' may perform a
stereotype of a tradition rather than a calibration to it, and
cultural-alignment evaluations of LLMs can be unstable across
elicitation choices. Three design features address this: each CRI
prompt is grounded in named contrastive-rhetoric sources rather than
improvised (Section~\ref{sec:scaling-panel}); agent behaviour is
checked against the corpus cohort corresponding to each lens; and all
runs are deterministic, removing sampling variation from the
elicitation.

\textbf{Internal validity.} Scores may be sensitive to prompt wording
and output format. All prompts are fixed before the experiments,
reproduced verbatim in Appendix~\ref{chap:appendix-prompts}, and
identical across cohorts, so wording effects cannot differentially
affect one cohort. Caching and fixed model versions eliminate drift
within an experiment; model version identifiers are recorded in
Appendix~\ref{chap:appendix-repro}.

\textbf{External validity.} ICNALE's topic control is a strength for
causal comparison and a limit on generality: two essay prompts in the
main cohorts, and a single prompt in the GRA validation set. Findings
therefore generalise to short argumentative learner essays of this
kind; extension to other genres, lengths, and scoring rubrics is
future work. The model set spans two closed and two open-weight
systems, but all are contemporary instruction-tuned LLMs; conclusions
are bounded to this model class.

\textbf{Statistical validity.} The GRA validation set is $n = 140$,
which limits the precision of E3's sensitivity and false-positive
estimates; all E3 statistics are reported with their $n$ and
confidence intervals, and the calibration/evaluation split is fixed by
seed before analysis. Across all experiments, confidence intervals are
the primary inferential device and no post-hoc significance testing is
performed (Section~\ref{sec:metrics}).

\section{Discussion}
% Interpret against H1-H4; scalability paradox; zero-sum boundary;
% relation to Dorleon & Shujaa's formal framework; practical
% implications for institutions; honest treatment of Part A limitations.
%% TODO: Week 9.
Lorem ipsum dolor sit amet, consectetur adipiscing elit.
